\chapter{Introducción}

En la actualidad, la toma de decisiones en grupo puede verse dificultada debido a la indecisión y la diversidad de opiniones. Situaciones cotidianas como elegir una película, planificar una salida cultural o decidir dónde cenar suelen requerir un tiempo significativo, llegando incluso a generar desacuerdos entre los participantes. \\

Este proyecto nace de la necesidad de optimizar los procesos de elección en entornos sociales. El objetivo principal es el desarrollo de una aplicación móvil orientada a la toma de decisiones en grupo de manera ágil e intuitiva. Mediante perfiles personalizables, los usuarios podrán crear núcleos sociales y proponer diversas alternativas que serán escogidas mediante distintos algoritmos de resolución, tales como la votación democrática, de forma aleatoria o por turnos.\\

\section{Motivación}
Este proyecto surge con el objetivo de reducir el tiempo invertido en decidir y mitigar los conflictos derivados de la indecisión y diversidad de opiniones. Abarcando desde elecciones triviales entre dos personas, como quién hace la comida, hasta decisiones más complejas con múltiples alternativas y un mayor número de participantes, como decidir qué película visualizar en común. Desde el punto de vista académico y técnico, esta aplicación constituye un reto significativo, al implicar la implementación completa de un sistema software de manera individual. Mientras que a lo largo del grado los proyectos de gran envergadura los hemos realizado en grupo, este proyecto me exige liderar todas las fases del ciclo de vida del software. Es especialmente relevante el reto técnico que supone el diseño y despliegue del back-end desde cero, expandiendo mis conocimientos más allá del front-end, mi área de especialización en el grado. Esta visión integral del desarrollo me ayudará a definir hacia dónde quiero orientar mi trayectoria profesional. Asimismo, la complejidad técnica y la aplicabilidad real del producto final constituyen un valor diferencial para mi currículum, alineándose con las competencias de autonomía y resolución de problemas complejos que demanda el sector tecnológico actual.

\begin{comment}
En la actualidad, la toma de decisiones en grupo puede verse dificultada debido a la indecisión y la diversidad de opiniones. Situaciones cotidianas como elegir una película, planificar una salida cultural o decidir dónde cenar suelen requerir un tiempo significativo, llegando incluso a generar desacuerdos entre los participantes. \\

El presente proyecto nace de esa necesidad de optimizar los procesos de elección en entornos sociales. El objetivo principal es el desarrollo de una aplicación móvil orientada a la toma de decisiones en grupo de manera ágil e intuitiva. Mediante perfiles personalizables, los usuarios podrán crear núcleos sociales y proponer diversas alternativas que serán escogidas mediante diversos algoritmos de resolución, tales como la votación democrática, de forma aleatoria o por turnos.


\section{Motivación personal}
La motivación de este proyecto surge con el objetivo de 
%agilizar y facilitar la toma de decisiones en el día a día en entornos sociales, 
mitigar la ineficiencia temporal y los conflictos interpersonales derivados de la indecisión grupal, abarcando desde elecciones triviales entre dos personas, como quién hace la comida, hasta decisiones más complejas con múltiples alternativas y un mayor número de participantes, como decidir qué película visualizar en común. El propósito último es favorecer un aprovechamiento más eficiente del tiempo compartido y disminuir los conflictos provocados por la indecisión y la diversidad de opiniones. \\

\section{Motivación académica y profesional}
Desde el punto de vista académico y técnico, esta aplicación constituye un reto significativo, al implicar la implementación completa de un sistema software de manera individual. Mientras que a lo largo del grado los proyectos de gran envergadura los hemos realizado en grupo, este proyecto me exige liderar todas las fases del ciclo de vida del software. Es especialmente relevante el reto técnico que supone el diseño y despliegue del back-end desde cero, expandiendo mis conocimientos más allá del front-end, mi área de especialización en el grado. Esta visión integral del desarrollo me ayudará a definir hacia dónde quiero orientar mi trayectoria profesional. Asimismo, la complejidad técnica y la aplicabilidad real del producto final constituyen un valor diferencial para mi currículum, alineándose con las competencias de autonomía y resolución de problemas complejos que demanda el sector tecnológico actual.
\end{comment}

\begin{comment}
Se pretende desarrollar una aplicación para dispositivos móviles de uso cómodo y sencillo, con la que organizar actividades de ocio en grupo como ver películas, ir al teatro o decidir dónde cenar. Los usuarios dispondrán de perfiles individualizados que podrán personalizar a su gusto, formarán parte de un grupo de amigos y podrán crear actividades en fechas concretas. Una vez creada la actividad, todos los miembros de un grupo podrán sugerir opciones propuestas por ellos mismos o sugeridas por la aplicación en base a sus preferencias. La decisión final se podrá tomar de varias formas: votación, aleatoriamente mediante una ruleta, por turnos entre los miembros, etc. \\


\section{Motivación personal}
La motivación de esta aplicación surge de manera personal, con el objetivo de agilizar y facilitar la toma de decisiones en mi día a día junto a mi grupo social, ayudando a resolver desde decisiones triviales entre dos personas, como quién hace la comida, hasta decisiones con múltiples opciones con un número elevado de personas, como decidir qué película visualizar en común. De esta manera, se reduciría en gran medida el tiempo invertido en decidir qué hacer y aprovecharíamos mejor el tiempo juntos. Además, disminuirían las discusiones provocadas por la indecisión y la diversidad de opiniones. \\

\section{Motivación académica y profesional}
En el ambito académico, esta aplicación presenta un gran reto, ya que es la primera vez que me enfrentaré al desarrollo de una aplicación completa de manera individual, dado que a lo largo de la carrera los proyectos de gran envergadura los hemos realizado en grupo. Me permitirá asimismo, crear el back-end desde cero, lo que supone una experiencia nueva, puesto que a lo largo de mi trayectoria estudiantil me he especializado más en el front-end. Me parece además, una gran oportunidad estar presente en todas las etapas del desarrollo software, lo que me permitirá definir más claramente hacía que tipo de puesto quiero orientar mi futuro laboral. Consiguientemente, se trata de una aplicación muy completa que puede impulsar notablemente mi currículum, debido a que las empresas valoran positivamente la experiencia y retos de este calibre.

\end{comment}

\section{Objetivos}

Se pretende desarrollar una aplicación multiplataforma con la siguiente funcionalidad:

\begin{itemize}
    \item Poder crear, editar y eliminar perfiles, grupos y actividades
    \item Integrar múltiples metodologías de toma de decisiones
    \item Poder guardar opciones en favoritos
    \item Disponer de opciones de personalización de la vista
    \item Garantizar la seguridad y privacidad de los datos personales
    \item Suministrar categorías de decisión por defecto
    \item Incorporar herramientas de filtrado
    \item Optimizar la toma de decisiones, reduciendo el tiempo y mitigando las discusiones

\begin{comment}
    \item Crear una aplicación accesible
    \item Crear una aplicación multiplataforma
    \item Tener una base de datos
    \item Que la aplicación disponga de tema claro y oscuro
    \item Que la aplicación este disponible en varios idiomas
    \item Tener varias formas distintas de toma de decisiones
    \item Perfiles personalizables
    \item Seguridad de los datos
    \item Poder crear grupos
    \item Poder editar grupos (miembros y cosas)
    \item Poder eliminar grupos
    \item Guardar amistades
    \item Guardar opciones favoritas
    \item Tener filtrado de opciones
    \item Disponer de categorias por defecto
    \item Tener un buen rendimiento
    \item Poder crear actividades
    \item Poder editar las actividades
    \item Poder editar los perfiles
    \item Poder editar las preferencias de visualización
    \item Guardar "ruletas" como favoritas
    \item Poder crearte un perfil
    \item Poder eliminar tu perfil
\end{comment}
    
\end{itemize}